\begin{theorem}[isCompletelyMultiplicative liouville]
\uses{liouville,IsCompletelyMultiplicative}
\label{isCompletelyMultiplicative_liouville}
\lean{ArithmeticFunction.isCompletelyMultiplicative_liouville}
% at /home/runner/work/trough/trough/library/PrimeNumberTheoremAnd/IwaniecKowalskiCh1.lean:1773.0-1783.7
\leanok
The Liouville function is completely multiplicative.
\end{theorem}
\begin{proof}
The Liouville function $\lambda(n)$ is defined as $(-1)^{\Omega(n)}$, where $\Omega(n)$ counts the total number of prime factors of $n$ with multiplicity. To show that $\lambda$ is completely multiplicative, we need to verify that $\lambda(1) = 1$ and that $\lambda(ab) = \lambda(a)\lambda(b)$ for all natural numbers $a$ and $b$.
\end{proof}
